\chapter*{Introduction générale}
\addcontentsline{toc}{chapter}{Introduction générale}
\thispagestyle{plain}

Juger si un fruit est «~prêt~» est une décision quotidienne, pour le consommateur comme pour le commerçant. Un fruit cueilli ou acheté trop tôt est dur et sans saveur~; un fruit laissé trop longtemps devient inutilisable et finit à la poubelle. Cette appréciation repose en grande partie sur la couleur et la texture de la peau, deux indices que la vision par ordinateur sait exploiter.

Les réseaux de neurones convolutifs (CNN) se sont imposés comme la méthode de référence pour reconnaître le contenu d’une image. Leur coût de calcul a longtemps réservé leur usage à des serveurs équipés de cartes graphiques. L’arrivée d’architectures conçues pour les appareils mobiles, comme \textbf{MobileNet}, change la donne~: un modèle de quelques mégaoctets peut désormais s’exécuter directement sur un téléphone, sans connexion et sans envoyer de photo à un serveur.

C’est dans ce cadre que s’inscrit notre projet du cours de \emph{Deep Learning}~: concevoir un système qui reconnaît, à partir d’une photographie prise avec un smartphone, \textbf{le fruit et son stade de maturité}. Nous avons retenu trois fruits (bananes, fraises, avocats) et onze classes, de «~non mûr~» à «~pourri~».

\begin{keybox}{Problématique}
\centering\itshape
Comment concevoir un classifieur de maturité de fruits à la fois suffisamment précis et suffisamment léger pour s’exécuter hors ligne sur un smartphone, et comment évaluer honnêtement sa fiabilité~?
\end{keybox}

Le projet couvre l’ensemble de la chaîne~: préparation de données annotées, entraînement par transfert d’apprentissage, évaluation rigoureuse, conversion au format TensorFlow Lite et intégration dans une application Android. Il vise aussi à expliquer les notions théoriques sous-jacentes (neurone, convolution, séparabilité en profondeur, transfert d’apprentissage), car leur compréhension conditionne les choix faits tout au long du travail.

\paragraph{Contributions.} Ce rapport présente~:
\begin{itemize}
  \item un jeu de données de 1\,598 images de fruits recadrés, réparties en 11 classes et découpées sans fuite d’information entre entraînement, validation et test~;
  \item un modèle MobileNetV2 affiné, atteignant 93,1\,\% d’exactitude sur un test indépendant, et une analyse détaillée de ses erreurs ;
  \item une série d’expériences négatives ou nuancées (callbacks mal configurés, augmentation plus forte, ajout de données d’une autre origine) qui éclairent les choix finaux~;
  \item une version quantifiée de 2,5\,Mo et une application Android complète, fonctionnant hors ligne, avec seuil de confiance et conseils par stade.
\end{itemize}

\paragraph{Organisation du rapport.} Le chapitre~\ref{chap:contexte} précise le contexte, les objectifs et la démarche. Les chapitres~\ref{chap:dl} et~\ref{chap:cnn} rappellent les bases du deep learning et des réseaux convolutifs~; le chapitre~\ref{chap:mobilenet} présente MobileNet et le transfert d’apprentissage. Le chapitre~\ref{chap:donnees} décrit les données, le chapitre~\ref{chap:modele} le modèle et son entraînement, et le chapitre~\ref{chap:resultats} les résultats. Les chapitres~\ref{chap:tflite} et~\ref{chap:android} traitent du déploiement et de l’application. Le chapitre~\ref{chap:discussion} discute les limites et les perspectives, avant la conclusion générale.
\clearpage
